% Volume IV: Law as an Epistemic Machine
% Part VII. Evidence Before Action
% Chapter 37: Admissibility
% Spine: proof-spines/ch037-spine.md

\chapter{Admissibility}
\label{ch:ch037}

Admissibility is treated here as a \textbf{type system} for institutional evidence. Material is not admitted merely because it exists or because it bears loosely on the facts; it must satisfy requirements of relevance, provenance, reliability, prejudice, and contestability. If \(\mathcal E\) is the space of all available material, the admissibility operator is
\[
\mathsf A_c:\mathcal E\to\{0,1,\bot\},
\]
\Definition\ where \(c\) is the procedural context, \(1\) means admitted, \(0\) refused, and \(\bot\) means unresolved pending further foundation.
The operator formalizes the thought that law does not ask only whether something is true. It asks whether that material is usable without deforming the path from appearance to institutionally warranted judgment. When those admissibility demands require forms of provenance or contestability that some injuries cannot realistically supply, the same rigor turns into an unmeetable burden that screens structurally important evidence out before it can be tested.
